%! AP Statistics — Unit 1, Lesson 1.7 — Warm-Up ANSWER KEY
\documentclass[10pt]{article}
\usepackage{apstats-article}
\usepackage{apstats-key}

%=============================================================================
\begin{document}

\pageheader{Unit 1, Lesson 1.7}{Warm-Up --- Answer Key}
\begin{tabularx}{\linewidth}{X X X}
Name:\;\ans{Key} & Date:\; & Period:\;
\end{tabularx}

\vspace{0.16in}

\noindent In Lesson 1.6 you described distributions using SOCS --- including
an \textit{informal} estimate of center. Today you will give those estimates
precise numerical values. Use the dotplot below.

\vspace{0.14in}

\begin{center}
\begin{tikzpicture}[scale=1.0]
  \draw[thick] (0.5,0) -- (6.2,0);
  \foreach \x/\lbl in {0.60/5, 1.20/10, 1.80/15, 2.40/20, 3.00/25,
                        3.60/30, 4.20/35, 4.80/40, 5.40/45} {
    \draw (\x,-0.12) -- (\x,0.12);
    \node[below, font=\small] at (\x,-0.12) {\lbl};
  }
  \node[below, font=\small\itshape] at (3.0,-0.56)
        {One-way commute time (minutes)};
  \filldraw[navy] (0.96, 0.38) circle (0.13);
  \filldraw[navy] (1.44, 0.38) circle (0.13);
  \filldraw[navy] (1.80, 0.38) circle (0.13);
  \filldraw[navy] (1.80, 0.76) circle (0.13);
  \filldraw[navy] (2.16, 0.38) circle (0.13);
  \filldraw[navy] (2.40, 0.38) circle (0.13);
  \filldraw[navy] (2.64, 0.38) circle (0.13);
  \filldraw[navy] (3.00, 0.38) circle (0.13);
  \filldraw[navy] (3.36, 0.38) circle (0.13);
  \filldraw[navy] (5.40, 0.38) circle (0.13);
\end{tikzpicture}
\end{center}

\vspace{0.16in}

\begin{enumerate}[leftmargin=1.5em, itemsep=8pt]

  \item How many students are represented in this dotplot? What is the
        variable being measured?\\[4pt]
        \ans{$n = 10$ students. The variable is one-way commute time, measured
        in minutes. It is quantitative.}

  \item What is the \textbf{range} of commute times? Show your calculation.\\[4pt]
        \ans{Range $= 45 - 8 = \mathbf{37}$ minutes.}

  \item List the data values in order from least to greatest. Then identify
        the \textbf{middle value} (or middle two values) of the sorted list.\\[4pt]
        \ans{Sorted: 8, 12, 15, 15, 18, 20, 22, 25, 28, 45.
        $n = 10$ (even), so the two middle values are the 5th and 6th: \textbf{18 and 20}.
        The median $= (18+20)/2 = \mathbf{19}$ minutes.}

  \item One commute time appears to be much larger than the rest. Which value
        is it, and how does it stand out in the dotplot?\\[4pt]
        \ans{The value 45 minutes stands out --- it is isolated far to the right
        of the other nine values, which cluster between 8 and 28 minutes. There
        is a large gap between 28 and 45.}

\end{enumerate}

\vspace{0.14in}

\begin{tcolorbox}[colback=greenbg, colframe=greenacc, boxrule=0.8pt, arc=2mm,
    left=4mm, right=4mm, top=3mm, bottom=3mm]
\small
\begin{itemize}[leftmargin=1.3em, itemsep=5pt]
  \item The \textbf{median} is exactly that middle value you identified above ---
        the precise measure of center that resists extreme values.
  \item Today we also calculate the \textbf{mean}, \textbf{IQR}, and
        \textbf{standard deviation}, and learn when each is the right choice.
  \item The large value at 45 minutes will motivate the \textbf{outlier rule}.
\end{itemize}
\end{tcolorbox}

\end{document}
